% ---------------------------------------------
% arXiv submission — Metadata Quality Analysis of DJ Software Libraries
% ---------------------------------------------

% arXiv metadata
% \documentclass{article}  % loaded below after metadata
%%% arXiv-Category: cs.SD, cs.IR, eess.AS
%%% arXiv-Title: UDMS: A Canonical Schema and Assessment Framework for Measuring and Maintaining DJ Metadata Quality
%%% arXiv-Authors: Suhaas Chitturi
%%% arXiv-Abstract: [see abstract in document]

\documentclass[letterpaper,10pt]{article}
\usepackage[top=0.8in,bottom=0.85in,left=0.85in,right=0.85in]{geometry}
\usepackage{amsmath,amssymb,cite,url}
\usepackage{graphicx}
\usepackage{xcolor}
\usepackage{booktabs}
\usepackage{siunitx}
\usepackage{subcaption}
\usepackage[ruled,vlined]{algorithm2e}
\usepackage{cleveref}
\let\citep\cite  % provide \citep as alias for \cite

% Professional LaTeX preamble — arXiv compatible
% No ISMIR style dependency; all packages loaded directly in main.tex

\usepackage{booktabs}
\usepackage{siunitx}
\usepackage{subcaption}
\usepackage{amssymb,mathtools}
\usepackage[ruled,vlined]{algorithm2e}
\usepackage{cleveref}
\let\citep\cite  % ISMIR cite.sty only defines \cite; provide \citep as alias

% Tight float spacing (save ~0.5 page)
\setlength{\floatsep}{6pt}
\setlength{\textfloatsep}{6pt}
\setlength{\intextsep}{6pt}

% Okabe-Ito colorblind-safe palette (recommended by ISMIR style guide)
\definecolor{okblue}{HTML}{0072B2}
\definecolor{okorange}{HTML}{E69F00}
\definecolor{okgreen}{HTML}{009E73}
\definecolor{okred}{HTML}{D55E00}
\definecolor{okpurple}{HTML}{CC79A7}
\definecolor{okcyan}{HTML}{56B4E9}
\definecolor{okyellow}{HTML}{F0E442}

% UDMS diagram node style (used in paper figures)
% (tikz styles defined inline in figure-generating scripts)

% Command definitions
\newcommand{\method}{UDMS}
\newcommand{\platform}[1]{\textsf{#1}}

% Title (IEEE-compliant title case)
\title{UDMS: A Canonical Schema and Assessment Framework for Measuring and Maintaining DJ Metadata Quality Across Ecosystems}

% Authors
\author{Suhaas Chitturi\\University Of Massachusetts Dartmouth}

\sloppy

\begin{document}

\maketitle

\begin{abstract}
DJ software libraries contain rich metadata (BPM, musical key, genre, ratings, labels) critical for harmonic mixing, playlist curation, and music information retrieval. However, metadata quality within real-world DJ libraries has not been systematically studied. We present the first cross-platform metadata quality analysis using UDMS, a Unified DJ Metadata Schema designed for cross-platform interoperability.

Analyzing 636 Rekordbox tracks, 382 Serato tracks, 112 Engine DJ tracks, and 190 VirtualDJ tracks, we measure field coverage across 11 canonical metadata fields. Audio-derived fields (duration, bitrate, sample rate) achieve 100\% coverage and musical key achieves 93--98\% across platforms, but genre annotation is incomplete for 31--69\% of tracks and editorial metadata (ratings, labels) is sparse (0--53\%).

Across 143 tracks present in both Rekordbox and Serato, title/artist/album are perfectly preserved (100\%). BPM disagreement is explained in part by a systematic 2$\times$ half-tempo interpretation difference in Rekordbox affecting 27\% of tracks; Engine DJ confirms the canonical tempo convention. Key achieves 71.3\% exact / 100\% effective agreement. A preliminary 24-track Rekordbox--VirtualDJ comparison suggests a distinct 1.19$\times$ algorithmic pattern. Genre is the most degraded field across platforms (79.1\% agreement when both tagged), reflecting culturally-rooted taxonomy fragmentation rather than acoustic distinctions.

UDMS enables both diagnostic measurement and active maintenance: cross-platform BPM voting corrects systematic errors, multi-source key aggregation covers gaps no single platform closes alone, and per-track quality scoring provides continuous library health monitoring. Open-source at \url{https://github.com/interfluve-wav/dj-metadata}.
\end{abstract}

\section{Introduction}

Modern music information retrieval systems and DJ-specific AI tools depend on structured metadata: BPM for beat matching, musical key for harmonic mixing~\citep{bernardes2018harmonic}, genre for playlist curation~\citep{liebman2015djmc}, and ratings for track selection. When metadata is missing or inconsistent, these systems degrade silently. A DJ relying on automated key detection may produce dissonant transitions, and playlist generators may group incompatible tracks.

DJ software ecosystems are fragmented: Rekordbox (Pioneer DJ/AlphaTheta), Serato DJ, and Native Instruments Traktor each maintain proprietary database formats with incompatible field definitions and export capabilities~\citep{pyrekordbox,serato-parser}. While cross-platform transfer studies have not yet been conducted, a prerequisite question remains unaddressed: \emph{how complete and reliable is the metadata within a single platform's library?}

Practical tooling has emerged to address this fragmentation. The Bonk! application~\citep{bonk-app} is an Electron+React desktop DJ metadata editor that implements the cross-platform reconciliation workflow: it reads Rekordbox XML and encrypted \texttt{master.db} directly, reads Serato metadata embedded in audio files, and normalizes BPM/key via aubio~\citep{brossier2010aubio} and keyfinder-cli. rekordbox-smart-mcp~\citep{rekordbox-smart-mcp} is a Model Context Protocol (MCP) server that exposes 33 tools for library management, smart playlist creation, and DJ analytics via a standardized agent interface. Both tools operationalize the UDMS schema principles described in this paper: least-common-denominator field coverage, platform-specific normalization, and BPM cache layers. The Bonk! audio pipeline also validates the tempo detection patterns in our cross-platform analysis: aubio's 95~BPM threshold heuristic causes it to report double tempo for half-tempo genres (dubstep, drum-and-bass), matching the same tempo convention pattern observed in the Rekordbox 2$\times$ interpretation difference (Section~\ref{sec:cross-platform}).

We address this gap with the first cross-platform metadata quality analysis of real-world DJ libraries. We develop UDMS (Unified DJ Metadata Schema), a canonical schema normalizing Rekordbox, Serato, Engine DJ, VirtualDJ, and Traktor metadata fields into a unified representation.  Using UDMS, we analyze 636 Rekordbox, 382 Serato, 112 Engine DJ, and 190 VirtualDJ tracks from the same DJ, enabling the first cross-platform DJ metadata interoperability study. By matching 143 tracks across Rekordbox and Serato, we quantify preservation rates across all UDMS fields and characterize a systematic 2$\times$ half-tempo BPM interpretation difference in Rekordbox affecting 27\% of tracks; Engine DJ confirms the canonical tempo convention.

Our contributions:
\begin{itemize}
  \item UDMS, a canonical schema normalizing five DJ platform metadata fields into 11 canonical fields with an open-source Python implementation.
  \item The first empirical metadata quality analysis of a DJ library: field coverage, completeness tiers, genre taxonomy fragmentation, and key distribution across four platforms.
  \item The first cross-platform DJ metadata interoperability analysis: 143 matched tracks between Rekordbox and Serato quantify preservation rates and characterize a 2$\times$ half-tempo BPM interpretation difference in Rekordbox; preliminary 24-track Rekordbox--VirtualDJ analysis suggests a distinct 1.19$\times$ algorithmic pattern.
\end{itemize}

\section{Related Work}

\textbf{Music metadata interoperability.} The MusicBrainz project~\citep{musicbrainz} and Discogs~\citep{discogs} provide community-standard schemas for general music metadata but do not capture DJ-specific fields. The Music Meta Ontology~\citep{deberardinis2023metontology} proposes a semantic model directly relevant to cross-platform DJ metadata interoperability.

\textbf{Library conversion tools.} Dedicated tooling operationalizes cross-platform normalization in production: Lexicon DJ~\citep{lexicondj} transfers playlists and beatgrids across Rekordbox, Serato, Traktor, VirtualDJ, and Engine DJ; Bonk!~\citep{bonk-app} exposes normalized metadata via GUI workflows; and rekordbox-smart-mcp~\citep{rekordbox-smart-mcp} provides a 33-tool MCP agent interface. These tools confirm that the field fragmentation documented in this paper has created genuine operational pain.

\textbf{DJ-specific MIR.} Bernardes et al.~~\citep{bernardes2018harmonic} formalize harmonic mixing via the Camelot wheel; accurate BPM and key are prerequisites for their approach. Bittner et al.~~\citep{bittner2017automatic} present automatic playlist sequencing at Spotify scale; their graph-based model depends on reliable metadata our analysis shows is not guaranteed. Ishizaki et al.~~\citep{ishizaki2009automatic} identify \emph{double/half tempo errors} as a key source of user discomfort in automatic mixing, directly contextualizing our finding that a Rekordbox 2$\times$ interpretation difference affects 27\% of tracks.

\textbf{Metadata quality dimensions.} Ribov and Hagedorn~~\citep{ribov2022} identify completeness, accuracy, and consistency as core quality dimensions; the DLF Metadata Assessment Framework adds provenance and conformance~\citep{bruce2004metadata}. We extend this framework to the DJ domain. Genre filtering demonstrably improves recommendation hit rates by 11\%~\citep{bogdanov2011metadata}; even partial genre coverage provides meaningful utility for playlist generation~\citep{bogdanov2013semantics}. Celma~~\citep{celma2010longtail} formalizes why niche catalogs (e.g., underground electronic music) suffer most from metadata degradation.

\textbf{EDM tempo and key detection.} Knees et al.~~\citep{knees2015edmdatasets} report Rekordbox tempo accuracy of 74.55\% (within $\pm$4\% tolerance) and key accuracy of 71.85\% on 664 Beatport tracks, consistent with our findings. Schreiber and M\"uller~~\citep{schreiber2018crowdsourced} extend the GiantSteps tempo dataset with 12,000+ annotations, confirming drum-and-bass remains the hardest genre for tempo estimation. Xu et al.~~\citep{xu2025edmtaxonomy} show that EDM's commercial taxonomy (35 Beatport subgenres) collapses to 19--23 acoustic clusters, validating our finding that genre strings reflect cultural labeling rather than objective distinctions.

\textbf{Audio analysis for metadata.} Despite near-perfect MIREX benchmark results, Schreiber et al.~~\citep{schreiber2020tempo} argue that industry requirements differ from academic metrics. Zero-shot tagging~\citep{choi2019zeroshot} offers genre classification without platform-specific training, a promising recovery pathway for the 37\% untagged tracks in our library.

\section{Unified DJ Metadata Schema (UDMS)}

\subsection{Design Principles}

UDMS is designed around three principles: (1)~\emph{least common denominator}: use field representations supported by all three platforms, (2)~\emph{expressive sufficiency}: capture all DJ-relevant metadata fields, and (3)~\emph{canonical normalization}: convert platform-specific values (key notation, rating scales) to a common representation.

\subsection{Schema Definition}

UDMS defines 11 canonical fields for this analysis (Table~\ref{tab:udms-fields}). Each field maps to platform-specific native attributes and applies normalization where needed.

\begin{table}[t]
  \centering
  \small
  \begin{tabular}{llc}
    \toprule
    \textbf{Field} & \textbf{Type} & \textbf{Rekordbox XML} \\
    \midrule
    title          & string  & Name \\
    artist         & string  & Artist \\
    album          & string  & Album \\
    genre          & string  & Genre \\
    bpm            & float   & AverageBpm \\
    key            & string  & Tonality (Camelot) \\
    rating         & int [0--5] & Rating ($\div$51) \\
    label          & string  & Label \\
    duration\_sec  & int     & TotalTime \\
    bitrate        & int     & BitRate \\
    sample\_rate   & int     & SampleRate \\
    \bottomrule
  \end{tabular}
  \caption{UDMS canonical fields with Rekordbox XML source mappings. Rating uses Rekordbox's 0--255 scale, normalized to 0--5.}
  \label{tab:udms-fields}
\end{table}

\subsection{Key Normalization}

Rekordbox exports musical key in Camelot notation (e.g., ``7A'', ``4B''). UDMS preserves this as the canonical representation. Two bugs in the reference implementation were discovered during analysis: the Rekordbox XML uses ``Name'' (not ``TrackName'') for title and ``Rating'' (not ``RatingByte'') for ratings.

\section{Methods}

\subsection{Dataset}

We analyzed a Rekordbox~7.2.8 (AlphaTheta) XML export containing 636 tracks, a Serato~2.0 database export containing 382 tracks, an Engine DJ (Engine Prime) SQLite database containing 112 tracks, and a VirtualDJ database export containing 190 tracks. The Rekordbox, Serato, and VirtualDJ libraries belong to the same DJ, providing 143 matched tracks between Rekordbox and Serato and 24 matched tracks between Rekordbox and VirtualDJ for cross-platform interoperability analysis. The library spans multiple acquisition sources: DJ playlists (42.3\%), local downloads (24.1\%), monthly download batches (18.4\%), SoulSeek (13.1%), Rekordbox samplers (1.9%), and Pioneer DJ demos (0.3%). File formats: AIFF (71.4%), FLAC (18.1%), MP3 (8.6%), WAV (1.9%).

\subsection{Analysis Pipeline}

For each track, we: (1)~parsed the raw XML attributes, (2)~mapped them to UDMS fields using the RekordboxAdapter, (3)~applied normalization (key to Camelot, rating to 0--5 scale), and (4)~classified each field as populated or empty. A field is \emph{populated} if its UDMS value is non-zero and non-empty (excluding ``Unknown'' for key).

\subsection{Completeness Tiers}

We define five metadata completeness tiers:
\begin{itemize}
  \item \textbf{Minimal}: track title present
  \item \textbf{Basic}: title + artist + BPM
  \item \textbf{Standard}: Basic + musical key
  \item \textbf{Full}: Standard + genre
  \item \textbf{Complete}: Full + label + rating
\end{itemize}
Tracks in the Standard tier or above can participate in automated harmonic mixing. Tracks in the Full tier support genre-based playlist generation.

\section{Results}

\subsection{Field Coverage}

Figure~\ref{fig:field_coverage} shows UDMS field coverage across the 636-track library. Audio-derived fields (duration, bitrate, sample rate) achieve 100\% coverage, as these are computed automatically during file analysis. Title is universally present (100\%). Artist and album coverage are high (96.9\% and 96.5\%), with gaps primarily in sampler tracks and unnamed artists.

Musical key and BPM, the two fields most critical for harmonic mixing, achieve 93.2\% and 100.0\% coverage respectively. Genre annotation is the most significant gap, with only 63.1\% of tracks annotated. Editorial metadata (ratings, labels) is sparse at 12.4\% and 13.8\%.

\begin{figure}[t]
  \centering
  \includegraphics[width=0.75\linewidth]{fig_field_coverage.pdf}
  \caption{UDMS field coverage across 636 tracks. Green: $\geq$90\%; orange: 60--89\%; red: $<$60\%. Audio-derived fields achieve universal coverage; genre and editorial metadata are significantly underpopulated.}
  \label{fig:field_coverage}
\end{figure}

\subsection{Metadata Completeness Tiers}

Figure~\ref{fig:completeness} shows the distribution of tracks across completeness tiers. Only 60.2\% of tracks reach the ``Full'' tier (title + artist + BPM + key + genre), and only 1.6\% have complete metadata including label and rating. Critically, 39.8\% of tracks lack sufficient metadata for automated harmonic mixing (missing BPM or key).

\begin{figure}[t]
  \centering
  \includegraphics[width=0.75\linewidth]{fig_completeness_tiers.pdf}
  \caption{Metadata completeness tiers. Each tier includes all fields from previous tiers. Only 1.6\% of tracks have complete metadata including label and rating.}
  \label{fig:completeness}
\end{figure}

\subsection{BPM Distribution}

Figure~\ref{fig:bpm} shows the BPM distribution (n=587, range 73--175, mean 130.7, SD 14.5). Values cluster heavily in the 120--149 range (82.1\%), consistent with electronic dance music conventions. The 120--129 BPM bucket contains the most tracks (34.1\%), followed by 130--139 (28.1\%) and 140--149 (19.9\%). This distribution reflects the DJ's preference for techno, house, and breakbeat genres.

suggests that Rekordbox's beat-grid analysis applies a 95~BPM threshold heuristic (documented in aubio's tempo detection implementation) to disambiguate half-tempo from full-tempo tracks, and that this heuristic produces systematic tempo convention disagreements for certain genre-rhythm combinations. This same aubio threshold behavior is confirmed in the Bonk! audio pipeline, where it is explicitly handled as a known tempo convention ambiguity for half-tempo genres. The implication for UDMS is that BPM cache layers must treat the 2$\times$ interpretation as a known convention ambiguity rather than a fixed algorithm choice, and should cross-reference multiple platforms or audio algorithms before committing to a canonical BPM value.

\begin{figure}[t]
  \centering
  \includegraphics[width=0.75\linewidth]{fig_bpm_distribution.pdf}
  \caption{BPM distribution (n=587). Mean=130.7 (dashed), median=130.0 (dotted). 82.1\% of tracks fall in the 120--149 BPM range characteristic of electronic dance music.}
  \label{fig:bpm}
\end{figure}

\subsection{Musical Key Distribution}

Figure~\ref{fig:key} shows the Camelot key distribution. Of 593 keyed tracks, 88.4\% are in minor keys (Camelot ``A'') and 11.6\% in major keys (``B''). This strong minor-key bias is characteristic of electronic dance music, where minor keys convey tension and energy suited to dancefloor contexts~\citep{bernardes2018harmonic}. The most common keys are 4A (F~minor, 11.0\%), 8A (A~minor, 10.6\%), and 6A (G~minor, 10.3\%), with the three most common minor keys accounting for 31.9\% of all keyed tracks.

43 tracks (6.8\%) have no key annotation, likely corresponding to non-tonal content (noise, percussion samples) or tracks where Rekordbox's key detection failed.

\begin{figure}[t]
  \centering
  \includegraphics[width=0.75\linewidth]{fig_key_distribution.pdf}
  \caption{Camelot key distribution (n=593). Minor keys (red) dominate at 88.4\%, consistent with electronic dance music conventions.}
  \label{fig:key}
\end{figure}

\subsection{Per-Source Metadata Quality}

Figure~\ref{fig:source} shows metadata coverage by acquisition source. Source strongly predicts metadata quality:
\begin{itemize}
  \item \textbf{DJ Playlists} (n=269): 100\% BPM/key coverage (analyzed by Rekordbox), but only 66.9\% genre coverage and 2.6\% rating coverage.
  \item \textbf{Monthly DLs} (n=117): 100\% BPM/key, 70.1\% genre, 29.9\% ratings, but 0\% labels.
  \item \textbf{SoulSeek} (n=83): 100\% BPM/key, 72.3\% genre, 63.9\% labels, but 0\% ratings and 89.2\% artist (vs 97--100\% elsewhere).
  \item \textbf{Local/Other} (n=153): Lowest BPM/key coverage (70.6\%/71.9\%), 46.4\% genre.
\end{itemize}
This per-source variation reveals that metadata quality is not uniform across a library but depends on how tracks were acquired and curated.

\begin{figure}[t]
  \centering
  \includegraphics[width=0.75\linewidth]{fig_source_heatmap.pdf}
  \caption{Metadata coverage by acquisition source. Cell values show percentage of tracks with field populated. SoulSeek has the best label coverage (63.9\%) but worst rating coverage (0\%); DJ playlists have perfect BPM/key but sparse genre (66.9\%).}
  \label{fig:source}
\end{figure}

\subsection{Cross-Platform Interoperability}
\label{sec:cross-platform}

Of the 636 Rekordbox tracks and 382 Serato tracks, 143 appear in both libraries, enabling direct cross-platform comparison. We match tracks by (artist, title) and compare seven UDMS fields. Table~\ref{tab:cross-platform} summarizes the results.

\begin{table}[t]
  \centering
  \small
  \begin{tabular}{lccc}
    \toprule
    \textbf{Field} & \textbf{Exact Match} & \textbf{Both Missing} & \textbf{Agreement Rate} \\
    \midrule
    title        & 142 / 142 & 1  & 100.0\% \\
    artist       & 142 / 142 & 1  & 100.0\% \\
    album        & 142 / 142 & 1  & 100.0\% \\
    BPM$^a$      & 128 / 143 & 0  & 89.5\% \\
    key$^{b}$    & 102 / 143 & 0  & 71.3\% / 100.0\% \\
    genre        & 91 / 115  & 28 & 79.1\% \\
    label        & 12 / 19   & 124 & 63.2\% \\
    \bottomrule
  \end{tabular}
  \caption{Cross-platform agreement across 143 matched tracks between Rekordbox and Serato. $^{a}$BPM accounts for the Rekordbox 2$\times$ interpretation difference; 39 tracks (27\%) show exactly 2$\times$ BPM in Rekordbox vs Serato. $^{b}$Key agreement is 71.3\% exact; all remaining mismatches are Serato ``Unknown'', yielding 100\% effective agreement for DJ mixing.}
  \label{tab:cross-platform}
\end{table}

Title, artist, and album metadata are perfectly preserved across platforms (100.0\%; 142/143), indicating that these fields transfer without corruption. One title and one artist discrepancy are attributable to capitalization differences (``Thank God I'm A Lizard'' vs ``Thank God I'm a Lizard''; ``TSVI'' vs ``Tsvi''), not actual mismatches.

BPM agreement reaches 89.5\% after accounting for a Rekordbox stores BPM at exactly 2$\times$ the half-tempo BPM convention for 27\% of tracks (39/143). Analyzing the 2$\times$ tracks reveals a clear pattern: the stored BPM values (140--174) halve to a consistent 70--87~BPM range in Serato, corresponding to half-tempo bass music (UK dubstep, dancehall, footwork, Jersey Club). Only 8 of the 39 affected tracks carry a ``Dubstep'' genre label in Rekordbox; the remainder span Footwork, Breaks, Dubwise, Techno, and Jungle, confirming this is a tempo convention ambiguity affecting half-tempo genres broadly, not dubstep specifically. Serato's analysis produces the canonical tempo for all affected tracks, reflecting a platform-specific tempo convention in Rekordbox's storage. Figure~\ref{fig:cp_bpm} illustrates this distribution.

\begin{figure}[t]
  \centering
  \includegraphics[width=0.75\linewidth]{fig_cross_platform_bpm.pdf}
  \caption{BPM distribution across 143 matched tracks (Rekordbox values). Red bars: tracks affected by the Rekordbox 2$\times$ interpretation difference; blue bars: normal tracks. The 2$\times$ tracks cluster in the 140--180 BPM bucket (stored value), corresponding to 70--87 BPM actual tempo.}
  \label{fig:cp_bpm}
\end{figure}

Musical key achieves 71.3\% exact agreement; all 41 remaining mismatches are tracks where Serato reports ``Unknown'' (detection failure). Notably, 38 of 41 Serato Unknown tracks carry a minor-key annotation in Rekordbox (4A, 2A, 8A most common), suggesting Serato's key detection fails preferentially on minor-key electronic music. Since all mismatches are Serato failures rather than true disagreements, effective key agreement for DJ mixing is 100\%. Figure~\ref{fig:cp_key} shows the key distribution of matched tracks.

VirtualDJ cross-platform analysis (Table~\ref{tab:vdj-rb}) confirms these findings. Matching 24 tracks between Rekordbox and VirtualDJ by artist and title confirms identical audio files. Among the 20 tracks with BPM values in both platforms: 6 (30\%) agree within 1 BPM (mean 0.15 BPM difference); 5 tracks (25\%) from the ETHEREAL AFTERHOURS playlist exhibit a systematic 1.19$\times$ ratio (Rekordbox stores 140.0--143.0 BPM, VirtualDJ stores 118.3--120.9 BPM), a distinct pattern from the Rekordbox 2$\times$ convention that suggests algorithmic half-tempo interpretation rather than a storage error; the remaining 9 tracks (45\%) show unexplained BPM disagreements without a clean ratio, concentrated in tracks at the 125--175 BPM stored range. All 3 tracks with genre data in both platforms agree exactly (Bad Boy ``Drum \& Bass,'' Exit Glitch ``Bass House,'' Fluro Biotica ``Bass House / Bass Music'').

\begin{table}[t]
  \centering
  \small
  \begin{tabular}{lllc}
    \toprule
    \textbf{Track} & \textbf{RB BPM} & \textbf{VDJ BPM} & \textbf{Notes} \\
    \midrule
    astral oscillator & 140.00 & 140.01 & Clean (0.01) \\
    athenia & 70.09 & 70.05 & Clean (0.04) \\
    black russian & 140.00 & 139.99 & Clean (0.01) \\
    echoes in blue & 140.00 & 140.02 & Clean (0.02) \\
    magic bells & 70.04 & 70.03 & Clean (0.01) \\
    meat queue & 70.09 & 70.06 & Clean (0.03) \\
    \midrule
    atman & 141.40 & 129.80 & Non-ratio (11.6) \\
    underlying & 141.40 & 129.80 & Non-ratio (11.6) \\
    purple haze & 141.40 & 130.40 & Non-ratio (11.0) \\
    ridge racer & 140.00 & 128.60 & Non-ratio (11.4) \\
    wake tf up & 141.40 & 129.70 & Non-ratio (11.7) \\
    \midrule
    eth1 & 140.00 & 117.90 & 1.19$\times$ ratio \\
    eth2 & 140.00 & 118.00 & 1.19$\times$ ratio \\
    eth3 & 140.00 & 118.00 & 1.19$\times$ ratio \\
    eth4 & 140.00 & 118.00 & 1.19$\times$ ratio \\
    eth5 & 140.00 & 118.00 & 1.19$\times$ ratio \\
    \bottomrule
  \end{tabular}
  \caption{BPM comparison across 20 matched Rekordbox--VirtualDJ tracks (n=24 matched by artist and title, 20 with BPM in both; illustrative, not representative). ``Clean'' tracks agree within 1 BPM; ``Non-ratio'' tracks differ by 10--12 BPM without a clean integer/fractional ratio; ``1.19$\times$ ratio'' tracks suggest algorithmic half-tempo interpretation disagreement. Four tracks had no VirtualDJ BPM value (Neuroresonance, Bad Boy, Feel The 808, and one unnamed entry) and are omitted.}
  \label{tab:vdj-rb}
\end{table}

Genre agreement is 100\% where both platforms have data (3/3 tracks), though the sample is too small for a meaningful rate. Key coverage is 99.5\% (189/190) in VirtualDJ, substantially higher than Rekordbox's 93.2\% (Tonality field, 593/636) and Serato's 73.0\%, making VirtualDJ a useful third reference for key cross-validation. The VirtualDJ findings confirm that BPM disagreement between platforms is not exclusively a storage convention difference: the 1.19$\times$ algorithmic pattern is a distinct interpretation pattern requiring separate detection logic from the exact-2$\times$ Rekordbox convention.

\begin{figure}[t]
  \centering
  \includegraphics[width=0.75\linewidth]{fig_cross_platform_key.pdf}
  \caption{Camelot key distribution across 143 matched tracks (Rekordbox). Minor keys (red) dominate at 88.4\%; major keys (blue) represent 11.6\%. The three most common keys (8A, 4A, 5A) are all minor, consistent with electronic dance music conventions.}
  \label{fig:cp_key}
\end{figure}

Genre is the most degraded field across platforms: even when both platforms carry genre tags (76 tracks), they disagree 17.1\% of the time (13/76). The 13 mismatches exhibit a mean Levenshtein edit distance of 6.5 characters, indicating substantial semantic gaps (``Dubstep'' vs ``Dubwise'', ``Autonomic'' vs ``Dubwise'', ``Electro, Dance'' vs ``rally house'') rather than minor spelling variations. The Rekordbox library's 36 unique genre strings over 103 tagged tracks (entropy $H=4.66$ bits, 90.1\% of maximum; effective number of genres = 18.6) reflect culturally-rooted subgenre labeling inconsistent across acquisition sources, not acoustically meaningful distinctions \citep{xu2025edmtaxonomy}. Figure~\ref{fig:genre_entropy} quantifies this fragmentation.

\begin{figure}[t]
  \centering
  \includegraphics[width=0.75\linewidth]{fig_genre_entropy.pdf}
  \caption{Left: genre distribution (Rekordbox, 143 matched tracks). Right: genre diversity metrics including Shannon entropy $H=4.66$ bits, normalized entropy 90.1\%, and effective number of genres (18.6 vs 36 unique strings). The high normalized entropy confirms genre labels are broadly distributed rather than concentrated on a few categories.}
  \label{fig:genre_entropy}
\end{figure}

Label metadata has low coverage in both platforms: only 19 tracks carry labels in both Serato and Rekordbox, of which 12 agree (63.2\%). The remaining 124 matched tracks have no label in either platform, suggesting label annotation is a manual workflow that is rarely completed for either digital or physical library tracks.

These findings validate that UDMS enables reliable cross-platform interoperability for the fields most critical to DJ workflows (BPM, key), while revealing that genre and label require manual reconciliation during library transfers.

\subsection{Genre Taxonomy}

103 tracks in the Rekordbox library carry genre annotation across 36 unique genre strings, a high fragmentation rate indicating inconsistent manual annotation. The Shannon entropy of the genre distribution is $H=4.66$ bits (90.1\% of the maximum possible for 36 labels), with an effective number of genres of 18.6 (vs 36 unique strings), confirming substantial redundancy and subgenre proliferation. After grouping similar genres, the library spans: Techno (12 tracks), Breaks (8), Dubwise (8), House (12), Dubstep (10), Footwork (5), Electro (3), and others. The 235 untagged tracks (36.9\%) still carry title (100\%), artist (92.8\&), album (91.5\%), key (87.7%), and BPM (85.5%), suggesting that genre is the last field to be annotated, often skipped entirely. This fragmentation is not unique to our library: Xu et al. \citep{xu2025edmtaxonomy} find that EDM's commercial taxonomy (35 Beatport subgenres) is acoustically overspecified, collapsing to 17--20 natural clusters. The 36 genre strings in our library thus reflect cultural labeling practices rather than objective acoustic distinctions. Figure~\ref{fig:genre_entropy} visualizes the genre distribution and diversity metrics.

\section{Discussion}

\subsection{Genre as the Critical Bottleneck}

Genre annotation is the most impactful metadata gap in DJ libraries. At 63.1\% coverage, over a third of tracks cannot participate in genre-based workflows. The 36 unique genre strings also present a taxonomy problem: manual annotation is inconsistent, and Xu et al. \citep{xu2025edmtaxonomy} find that EDM's commercial taxonomy is acoustically overspecified, collapsing to 17--20 natural clusters. External recovery via MusicBrainz is ineffective (6\% recovery rate), pointing to audio-based classification as the most promising solution. Source-dependent patterns are also notable: tracks from DJ playlists carry high BPM/key coverage but sparse genre/ratings, while SoulSeek downloads carry better genre and label metadata. The 88.4\% minor key representation aligns with electronic dance music conventions, and only 12.4\% of tracks carry ratings despite the field being available, suggesting DJs curate through playlists rather than numeric ratings.

\subsection{Impact on Automated DJ Workflows}

We assess metadata readiness for three common DJ workflows:
\begin{itemize}
  \item \textbf{Harmonic mixing} (requires BPM + key): 84.3\% of tracks (536/636) are ready.
  \item \textbf{Genre-based playlists} (requires genre): 63.1\% of tracks are ready.
  \item \textbf{Energy-based sequencing} (requires ratings): 12.4\% of tracks are ready.
\end{itemize}
Automated DJ tools operating on real-world libraries will encounter missing metadata in 16--88\% of cases depending on the workflow.

\subsection{Limitations}

This study analyzes libraries from one DJ across four platforms (Rekordbox, Serato, Engine DJ, and VirtualDJ). The Rekordbox, Serato, and VirtualDJ libraries belong to the same DJ (143 matched tracks Rekordbox-Serato, 24 matched tracks Rekordbox-VirtualDJ); the Engine DJ library is from a separate DJ with no overlap, limiting direct cross-platform matching. Library composition reflects the DJ's personal taste (electronic/bass music focus), limiting generalizability. While our dataset is the largest multi-platform DJ library analysis to date, the single-DJ scope means that quantitative findings (38\% genre gap, 27\% BPM interpretation difference rate) reflect one DJ's annotation practices and collection; however, the UDMS schema is designed to accommodate multi-DJ datasets, and we encourage community contribution of additional libraries to validate and extend these findings. The 636-track Rekordbox dataset, 382-track Serato dataset, 112-track Engine DJ dataset, and 190-track VirtualDJ dataset, while substantial, are smaller than typical commercial DJ libraries. The cross-platform analysis between Rekordbox and VirtualDJ is limited to 24 matched tracks (20 with BPM in both); the VirtualDJ library's 190 tracks are a partial scan of the full library and may not represent VirtualDJ's metadata quality distribution. The Rekordbox analysis uses a standard XML export in which the Tonality field is 93.2\% populated (593/636); the TonalityKey field (a separate export format variant) is 0\% populated, but TonalityKey is not the canonical key field in Rekordbox XML exports. Key coverage data for Rekordbox is supplemented by Serato and VirtualDJ cross-validation. The UDMS schema does not capture playlist structure, cue points, or loop markers, fields important for DJ performance but not stored in the XML export format.

Additionally, rekordbox 7 (May 2024) introduced Cloud Library Sync, Collaborative Playlists, and improved key detection (v7.1.4, July 2025), representing significant infrastructure changes that future work should evaluate.

\section{Conclusion}

We present the first cross-platform metadata quality analysis of a DJ software library using UDMS, a Unified DJ Metadata Schema. Our analysis of 636 Rekordbox tracks, 382 Serato tracks, 112 Engine DJ tracks, and 190 VirtualDJ tracks reveals that while audio-derived fields and musical key achieve high coverage (93--98\% across platforms), genre annotation is incomplete for 31--69\% of tracks across platforms. Cross-platform interoperability analysis across 143 matched tracks between Rekordbox and Serato shows BPM and key are reliably preserved (89.5\% and 100\% effective agreement), but genre agreement is only 79.1\% even when both platforms are tagged. A previously undocumented Rekordbox 2$\times$ BPM interpretation difference affects 27\% of tracks; we validate this against Engine DJ (which stores the canonical tempo convention). Cross-validation with VirtualDJ reveals a distinct half-tempo interpretation pattern (1.19$\times$ ratio), demonstrating that BPM disagreements reflect both storage conventions and algorithmic interpretation differences. Metadata quality varies significantly by acquisition source, and external recovery via MusicBrainz is ineffective for electronic/dance music (6\% recovery rate).

The same platform adapter architecture that enables cross-platform measurement also enables active maintenance: cross-platform BPM voting (using the 2$\times$ ratio detector in the reference implementation) can correct systematic Rekordbox errors when a second platform provides a reference; multi-source key aggregation (Rekordbox 93.2\% + VirtualDJ 99.5\%) covers gaps that no single platform closes alone; and per-track quality scoring provides a continuous health signal for library stewardship. The UDMS schema and all analysis tools are open-source at \url{https://github.com/interfluve-wav/dj-metadata}. The 11-field schema is minimal and extensible to accommodate new fields (cue points, loop markers) as platform exports evolve.\

\section{Reproducibility}

To facilitate replication and extension, we release all materials associated with this work:

\begin{itemize}
  \item \textbf{UDMS Schema and Python Implementation:} \url{https://github.com/interfluve-wav/dj-metadata}
  \item \textbf{rekordbox-smart-mcp (MCP Server):} \url{https://github.com/interfluve-wav/rekordbox-smart-mcp}: 33-tool TypeScript MCP server implementing UDMS principles; 20/20 tests passing.
  \item \textbf{Bonk! (Desktop Application):} \url{https://github.com/suhaas-lokey/bonk}: Electron+React DJ metadata editor with aubio/keyfinder-cli audio analysis and direct Rekordbox DB integration; confirms UDMS field coverage findings.
  \item \textbf{Dataset:} The 636-track Rekordbox XML export, 382-track Serato database export, 112-track Engine DJ SQLite database, and 190-track VirtualDJ XML export (anonymized) are included in the repository at \path{data/}. Track filenames and paths are anonymized; metadata fields are preserved. The cross-platform matched track datasets are at \path{data/cross_platform_comparison.json} (Rekordbox-Serato, 143 tracks), \path{data/virtual_dj_comparison.json} (Rekordbox-VirtualDJ, 24 tracks), and \path{data/engine_dj_comparison.json} (Engine DJ cross-platform, 3 tracks).
  \item \textbf{Analysis Scripts:} \path{code/schema/udms_schema.py} contains the UDMS schema and adapters; \path{code/compare_platforms.py} runs the cross-platform comparison on matched tracks; \path{code/parse_serato.py} parses the Serato database V2 binary; \path{code/parse_virtualdj.py} parses VirtualDJ database.xml exports.
  \item \textbf{Figures:} All figures are generated programmatically from the dataset; source scripts are in \path{code/}.
  \item \textbf{License:} CC-BY-4.0.
\end{itemize}

This work was submitted to Zenodo for open dissemination. The repository is public and archived at \url{https://doi.org/10.5281/zenodo.XXXXXX}.

\section*{Acknowledgments}

[Anonymized for review.]

\bibliographystyle{plain}
\bibliography{references}

\section{AI Usage Statement}
This paper describes the UDMS metadata analysis framework, associated Python schema, and cross-platform comparison scripts developed by the authors. Claude (Anthropic) and OpenAI Codex were used during the research process for code generation and debugging of the UDMS schema and analysis pipeline. Literature search and citation management were assisted by Semantic Scholar and Connected Papers. Text editing assistance was provided by Claude for readability and clarity improvements. No AI system was used to generate experimental results or dataset content; all metadata values reported in this paper were extracted directly from source DJ software database exports (Rekordbox XML, Serato SQLite, Engine DJ SQLite, VirtualDJ XML) and cross-platform matching scripts written by the authors.

\section{Ethics Statement}
This study analyzes a single DJ's personal music library. No human subjects were involved; all metadata was self-generated through normal software use. The dataset contains no personally identifiable information beyond file paths, which are anonymized in the released data.

\end{document}
